/subsection{NVIDIA Tesla: A unified graphics and computing architecture}
This article explains how GPU-performance took a big leap with NVIDIAs Tesla-architecture by joining the processors of vertex- and pixel-processing to a single unified processor architecture design. By combining these programs to the same core, it enabled dynamic load management and therefore increased processor utilization compared to the fixed ratios. Also Tesla enabled general-purpose-computing with a GPU (GPGPU) by creating the powerful parallel-computing-platform called Compute Unified Device Architecture (CUDA) which is programmable in C.\\
The traditional pipeline with fixed ratios of pixel- and vertex processors is rarely balanced because the size of triangles are largely different and therefore some core sit still waiting for new tasks while the others are over-utilized. By combining these processes to a unified processor design (Streaming Multiprocessor) it not only allows for better load balancing but also lowers cost and the complexity of the chip-design by having less specialized hardware. Tesla also implemented the Single-Instruction-Multiple-Thread (SIMT)-model the first time. This clusters independent threads into warps which allows a more efficient execution of branched code.\\
In our opinion the Tesla architecture is a historic milestone of universal usage of GPUs for high-performance computing by providing more generalized core with the necessary interface (CUDA) to program them efficiently. These concepts are still in use in every modern NVIDIA-GPU and consolidates their marked dominance in AI and HPC.